\clearpage\fancyhead[L]{\small PALLAS.}\hypertarget{native-384-lower}{}\label{native-384-lower}
\begin{center}\textbf{PALLAS.}\par\vspace{4pt}\hrule\vspace{4pt}\fontsize{10}{12}\selectfont Göttingische gelehrte Anzeigen. No. 199. pp. 1985\;..\;1990. 12 December 1814.\par\vspace{4pt}\hrule\vspace{4pt}\end{center}

\par\indent Professor \textsc{Gauss}’s work on the intricate theory of the motions of Pallas, which we have mentioned several times in these pages, has again found fine confirmation in observations of the planet’s opposition in October of this year, the \emph{tenth} since its discovery. We give these observations here and the results \emph{immediately} derived from them, leaving the further implications for the theory to another occasion. In his prediction of the planet’s course for this period of visibility, made more than a year ago, Mr. \textsc{Nicolai}, adjunct at the Seeberg Observatory, used only the osculating elements of the ninth opposition (which was of course bound to be quite sufficient for finding the planet). Professor \textsc{Gauss} thought it of interest to determine, before the planet reappeared, how those elements would be corrected by Jupiter’s continuing influence from the ninth to the tenth opposition. Already at the beginning of the year he had Mr. \textsc{Möbius} carry out this considerable task under his direction. Mr. \textsc{Möbius} was then pursuing mathematical and astronomical studies here with exceptional zeal. The results of this calculation have already been printed in more detail in the astronomical yearbook for 1817, which has just appeared. We give here only the prediction of the opposition itself, for comparison with the actual observational results.

\begin{center}25 October 1814. $12^{\mathrm u}\;39^{\mathrm m}\;50^{\mathrm s}5$ Mean time in Göttingen\end{center}
\begin{center}\fontsize{10}{12}\selectfont\setlength{\tabcolsep}{3pt}\renewcommand{\arraystretch}{1.15}\begin{tabular}{lcr}
True longitude & $.\;.\;.\;.\;.\;.\;.\;.\;.\;.\;.\;.\;.$ & $31^\circ\;58'\;28''4$\\
Heliocentric latitude & $.\;.\;.\;.\;.\;.\;.\;.\;.\;.$ & $23\;38\;32.7\;\mathrm{S.}$\\
\end{tabular}\end{center}

\SourceLeaf{385}\Needspace{10\baselineskip}\par\indent The following observation was made with the ring micrometer at the observatory here.

\begin{center}\fontsize{9}{11}\selectfont\setlength{\tabcolsep}{3pt}\renewcommand{\arraystretch}{1.15}\begin{tabular}{c|ccc|ccc|ccc}
1814 & \multicolumn{3}{c|}{\shortstack{Mean time\\in Göttingen}} & \multicolumn{3}{c|}{\shortstack{Apparent right\\ascension}} & \multicolumn{3}{c}{\shortstack{Apparent\\declination}} \\
\hline
Sept.\ 16 & $12^{\mathrm u}$ & $1^{\mathrm m}$ & $1^{\mathrm s}$ & $46^{\circ}$ & $16'$ & $23''4$ & $11^{\circ}$ & $4'$ & $20''9$ S. \\
\end{tabular}\end{center}

\par\indent This confirmed the accuracy of the ephemeris. There could be no doubt that the planet could be observed with fixed instruments at other observatories, so there would have been no point in continuing the ring-micrometer observations. The observation above, moreover, was based on numerous comparisons and was marked as very good in the logbook.

\Needspace{19\baselineskip}\par\indent At the \emph{Seeberg} Observatory, the observers rightly confined themselves to observations with the transit telescope, since the planet was too faint to expect reliable declinations from the \textsc{Dollond} quadrant. The right ascensions observed by Mr. von \textsc{Lindenau} and Mr. \textsc{Nicolai} are as follows:

\begin{center}\fontsize{9}{11}\selectfont\setlength{\tabcolsep}{3pt}\renewcommand{\arraystretch}{1.15}\begin{tabular}{c|ccc|ccc}
1814 & \multicolumn{3}{c|}{\shortstack{Mean time\\at Seberg}} & \multicolumn{3}{c}{\shortstack{Apparent right\\ascension}} \\
\hline
Sept.\ 17  & $15^{\mathrm u}$ & $19^{\mathrm m}$ & $41^{\mathrm s}5$ & $46^{\circ}$ & $18'$ & $26''2$ \\
18         & 15 & 15 & 52.0 & 46 & 20 & 2.8 \\
19         & 15 & 12 & 1.0  & 46 & 21 & 17.0 \\
20         & 15 & 8  & 8.0  & 46 & 22 & 0.7 \\
Oct.\ 12 & 13 & 36 & 19.2 & 45 & 2  & 4.5 \\
13         & 13 & 31 & 52.0 & 44 & 54 & 14.6 \\
15         & 13 & 22 & 53.1 & 44 & 37 & 25.8 \\
21         & 12 & 55 & 28.0 & 43 & 39 & 51.6 \\
28         & 12 & 22 & 45.3 & 42 & 21 & 48.5 \\
\end{tabular}\end{center}

\Needspace{14\baselineskip}\par\indent The complete original observations made by Mr. \textsc{Schumacher} with the eight-foot mural quadrant at the \emph{Mannheim} Observatory were sent in and reduced here with the greatest care by Mr. \textsc{Encke}. The results were as follows:

\begin{center}\fontsize{9}{11}\selectfont\setlength{\tabcolsep}{3pt}\renewcommand{\arraystretch}{1.15}\begin{tabular}{c|ccc|ccc|ccc}
1814 & \multicolumn{3}{c|}{\shortstack{Mean time\\in Mannheim}} & \multicolumn{3}{c|}{\shortstack{Apparent right\\ascension}} & \multicolumn{3}{c}{\shortstack{Apparent\\declination}} \\
\hline
Oct.\ 12 & $13^{\mathrm u}$ & $36^{\mathrm m}$ & $17^{\mathrm s}4$ & $45^{\circ}$ & $2'$  & $0''3$  & $19^{\circ}$ & $12'$ & $38''4$ S. \\
15         & 13 & 22 & 51.5 & 44 & 37 & 23.7 & 20 & 6  & 54.6 \\
21         & 12 & 55 & 27.0 & 43 & 39 & 57.5 & 21 & 49 & 52.1 \\
22         & 12 & 50 & 48.3 & 43 & 29 & 15.3 & 22 & 6  & 3.8  \\
\end{tabular}\end{center}

\par\indent In the last two observations the planet was so faint that it could be made out only with difficulty.

\Needspace{18\baselineskip}\par\indent At the new \emph{Königsberg} Observatory (of whose great activity we have already given several examples in these pages), Professor \textsc{Bessel} made the following observations:

\begin{center}\fontsize{9}{11}\selectfont\setlength{\tabcolsep}{3pt}\renewcommand{\arraystretch}{1.15}\begin{tabular}{c|ccc|ccc|ccc}
1814 & \multicolumn{3}{c|}{\shortstack{Mean time\\in Königsberg}} & \multicolumn{3}{c|}{\shortstack{Apparent right\\ascension}} & \multicolumn{3}{c}{\shortstack{Apparent\\declination}} \\
\hline
Oct.\ 14 & $13^{\mathrm u}$ & $27^{\mathrm m}$ & $30^{\mathrm s}6$ & $44^{\circ}$ & $46'$ & $15''3$ & $19^{\circ}$ & $48'$ & $30''8$ S. \\
15         & 13 & 23 & 0.4  & 44 & 37 & 38.6 & 20 & 6  & 20.0 \\
18         & 13 & 9  & 22.9 & 44 & 10 & 8.3  & 20 & 58 & 58.9 \\
19         & 13 & 4  & 48.3 & 44 & 0  & 25.9 & 21 & 15 & 52.6 \\
Nov.\ 1    & 12 & 3  & 57.3 & 41 & 33 & 59.9 & 24 & 31 & 29.8 \\
2          & 11 & 59 & 12.6 & 41 & 21 & 46.2 & 24 & 44 & 9.7  \\
3          & 11 & 54 & 27.4 & 41 & 9  & 24.1 & 24 & 56 & $10\;6$ \\
4          & 11 & 49 & 42.0 & 40 & 56 & 59.6 & 25 & 7  & 41.8 \\
\end{tabular}\end{center}

\SourceLeaf{386}\par\indent This excellent observer also complains of the difficulties the planet’s faintness posed for the observations.

\Needspace{31\baselineskip}\par\indent The following table shows the results of comparing all these observations with the most recently determined elements. (We need only note that the epoch of the mean longitude used here differs by a few seconds, owing to a method of Professor \textsc{Gauss}’s own that is to be justified more fully elsewhere.)

\begin{center}
\fontsize{10}{12}\selectfont\setlength{\tabcolsep}{5pt}\renewcommand{\arraystretch}{1.15}\begin{tabular}{l|cc|l}
 & \multicolumn{2}{c|}{Difference} & \\
1814 & R.A. & Decl. & Observer \\
\hline
Sept.\ 16 & $-9''2$ & $+5''9$ & Gauss \\
17 & $-1.6$ & & von Lindenau \\
18 & $-6.2$ & & von Lindenau \\
19 & $-10.7$ & & von Lindenau \\
20 & $-7.5$ & & von Lindenau \\
Oct.\ 12 & $-4.9$ & & von Lindenau \\
12 & $-3.6$ & $+2.2$ & Schumacher \\
13 & $-8.8$ & & von Lindenau \\
14 & $-9.4$ & $+12.0$ & Bessel \\
15 & $-5.2$ & $+6.7$ & Bessel \\
15 & $-6.5$ & & von Lindenau \\
15 & $-7.7$ & $+5.7$ & Schumacher \\
18 & $-2.7$ & $+13.2$ & Bessel \\
19 & $-4.0$ & $+2.4$ & Bessel \\
21 & $-2.2$ & & von Lindenau \\
21 & $-12.0$ & $+4.5$ & Schumacher \\
22 & $-0.4$ & $-1.3$ & Schumacher \\
28 & $+1.8$ & & von Lindenau \\
Nov.\ 1 & $-0.4$ & $+8.5$ & Bessel \\
2 & $-5.1$ & $+10.6$ & Bessel \\
3 & $-5.9$ & $-2.8$ & Bessel \\
4 & $-7.9$ & $-21.9$ & Bessel \\
\end{tabular}
\end{center}

\Needspace{10\baselineskip}\par\indent Finally, the result obtained from these observations for the \emph{tenth opposition} itself is as follows:

\begin{center}25 October 1814. $12^{\mathrm u}\;33^{\mathrm m}\;22^{\mathrm s}$ Mean time in Göttingen\end{center}
\begin{center}\fontsize{10}{12}\selectfont\setlength{\tabcolsep}{3pt}\renewcommand{\arraystretch}{1.15}\begin{tabular}{lcr}
True longitude & $.\;.\;.\;.\;.\;.\;.\;.\;.\;.\;.\;.$ & $31^\circ\;58'\;11''3$\\
Geocentric latitude & $.\;.\;.\;.\;.\;.\;.\;.\;.\;.$ & $37\;20\;53.2\;\mathrm{S.}$\\
\end{tabular}\end{center}

\Needspace{21\baselineskip}\par\indent A precise comparison of this result with the most recently determined elements, duly corrected for Jupiter’s perturbations, gave such small differences that the basic elements require no further correction on this account. We add only the following overview of how closely all the oppositions observed so far agree with Mr. \textsc{Gauss}’s theory.

\begin{center}
\fontsize{10}{12}\selectfont\setlength{\tabcolsep}{5pt}\renewcommand{\arraystretch}{1.15}\begin{tabular}{ll|cc}
 & & \multicolumn{2}{c}{Difference} \\
Oppos. & & Mean longitude & Helioc. latitude \\
\hline
I & 1803 & $-1''3$ & $-5''3$ \\
II & 1804 & $-12.9$ & $-11.9$ \\
III & 1805 & $-20.4$ & $-0.4$ \\
IV & 1807 & $+2.5$ & $-6.8$ \\
V & 1808 & $+23.5$ & $+9.0$ \\
VI & 1809 & $+21.5$ & $-3.1$ \\
VII & 1811 & $+14.1$ & $-7.2$ \\
VIII & 1812 & $-16.1$ & $+4.0$ \\
IX & 1813 & $-10.9$ & $-0.0$ \\
X & 1814 & $-14.0$ & $+2.4$ \\
\end{tabular}
\end{center}
